\subsubsection{Structure des matériaux de type spinelle (\ce{AB2X4})}

\setcounter{equation}{0}
\setcounter{Q}{0}

Certains spinelles peuvent être utilisés dans le domaine des batteries. En effet, ce type de matériaux présente une structure en 3D avec des canaux pouvant accueillir des ions lithium.

\`{A} l'état solide, l'oxyde de lithium-manganèse \ce{Li_x Mn_2 O_4} cristallise dans un réseau de type spinelle, comme présenté ci-dessous : 

\begin{center}
\includegraphics[width=10cm]{Ex12/Fig01}
\captionof{figure}{Projection 2D de la structure \ce{LiMn2O4} (a) et structure 3D (b)}
\end{center}


\Q Rappeler la différence entre spinelle normal et spinelle inverse.
\solution{Les spinelles présentent deux fois plus de sites tétraédriques (T) qu'octaédriques (O), on a donc deux occupations possibles des sites : 
\begin{align*}
\ce{[A]_O [B_2 ]_T O4} &\implies \text{spinelle inverse} \\
\ce{[B]_O [A B ]_T O4} &\implies \text{spinelle normal} \\
\end{align*}}

\Q \`{A} partir d'un calcul du degré d'oxydation, montrer que \ce{LiMn2O4} est un composé de valence mixte.
\solution{
Par électroneutralité, on a : 
\begin{align*}
no ( \ce{Li} ) + 2 \cdot no ( \ce{Mn} ) + 4 \cdot no ( \ce{O} ) = 0
\end{align*}
Avec les électronégativités, on a $\chi ( \ce{Li} ) < \chi ( \ce{Mn} ) < \chi ( \ce{O} )$, on a donc $no ( \ce{O} ) = - II$ et $no ( \ce{Li} ) = + I$ : 
\begin{align*}
2 no ( \ce{Mn} ) = VII \implies no( \ce{Mn}) = III \quad no ( \ce{Mn} ) = IV
\end{align*}
Les degrés d'oxydation stable du manganèse permettent de déterminer la seule solution.}

Dans une structure spinelle, les ions \ce{O^{2-}} occupent les n{\oe}uds d'un réseau cubique à faces centrées, les ions métalliques se répartissant entre les sites octaédriques et tétraédriques du réseau.

\Q Indiquer la position des sites octaédriques et tétraédriques dans le réseau cubique à faces centrées. Dénombrer ces sites. 
\solution{On va dénombrer le nombre de sites octaédrique et tétraédrique, la multiplicité vaut : 
\begin{align*}
Z ( \ce{O^{2-}} ) &= 8 \times \frac{1}{8} +  6 \times \frac{1}{2} = 4 \\
Z ( Td ) &= 8 \\
Z ( Oh ) &= 12 \times \frac{1}{4} + 1 = 4
\end{align*}}

Dans ces structures, les ions manganèse au degré d'oxydation III et IV occupent les sites octaédriques. Cependant, durant des cycles de charge/décharge rapides ou menées à trop grande tension, les ions manganèse migrent parfois dans des sites tétraédriques.

\Q En utilisant le \og critère stérique \fg{}, justifier pourquoi l'ion manganèse peut facilement migrer d'un site octaédrique vers un site tétraédrique.
\solution{On peut calculer le rapport entre le rayon de l'anion et du cation : 
\begin{align*}
r( \ce{Mn^{3+}}, \text{HS}) &= 65 ~ \si{\pico\meter} &\implies \rho_1 = 0.46\\
r( \ce{Mn^{3+}}, \text{BS}) &= 58 ~ \si{\pico\meter} &\implies \rho_2 = 0.41\\
r( \ce{Mn^{4+}} ) &= 53 ~ \si{\pico\meter} &\implies \rho_3 = 0.38
\end{align*}
Selon le \og critère stérique \fg{}, les cations manganèse IV et III-BS (resp. III-HS) occupent les sites tétraédriques (resp. octaédriques). Toutes ces valeurs sont "proches" de la valeur de transition Td-Oh.}

\Q Donner la configuration électronique des ions manganèse (III) et (IV).
\solution{
D'après les règles de l'Aufbau, on a : 
\begin{align*}
[\ce{Mn}] &= ~ _{18}[\ce{Ar}] 4s^2 3d^5 \\
[\ce{Mn^{3+}}] &= ~ _{18}[\ce{Ar}] 4s^0 3d^4 \\
[\ce{Mn^{4+}}] &= ~ _{18}[\ce{Ar}] 4s^0 3d^3 \\
\end{align*}
}

\Q En utilisant le \og critère énergétique \fg{}, justifier pourquoi l'ion manganèse (III) peut facilement migrer d'un site octaédrique vers un site tétraédrique. On présumera que sa configuration est haut-spin (HS) et on la précisera.
\solution{
La configuration électronique de Mn(III) HS est  : 
\begin{align*}
(t_{2g})^3 (e_g)^1
\end{align*}
On va chercher à estimer l'énergie de stabilisation du champ cristallin (ESCC) dans les trois configurations possibles du manganèse (III) : 
\begin{align*}
ESCC (Oh) &= - \frac{3}{5} \Delta_O \\ 
ESCC (Td) &= - \frac{2}{5} \Delta_T = - \frac{8}{45} \Delta_O
\end{align*}
On peut alors calculer la différence d'ESCC :
\begin{align*}
ESCC ( Oh, \text{HS} ) - ESCC (Td) &= -\frac{3}{5} \Delta_O  + \frac{8}{45} \Delta_O \\
&= - \frac{19}{45} \Delta_O \approx - 0.422 \Delta_O
\end{align*}
On trouve une énergie de stabilisation inférieure à \SI{1}{\electronvolt}. Il semble tout à fait possible qu'à des tensions importantes, la structure cristallographique soit impactée.
}

Les atomes de lithium occupent préférentiellement des sites tétraédriques.

\Q En déduire si la structure est un spinelle normal ou inverse s'il n'y a pas d'altération de la charpente.
\solution{\`{A} priori, en l'absence de migration des ions manganèse, on a 
\begin{align*}
\ce{[Li]_T [Mn_2]_O O4}
\end{align*}
Il s'agit d'un spinelle inverse.
}

%\Q En utilisant le \og critère stérique \fg{}, pouvez-vous indiquer si la structure est déformée par l'insertion du lithium lors de la charge ? 
Du fait des migrations du manganèse dans la batterie, le composé \ce{Li_x Mn_2 O_4} a une durée vie relativement faible. Lors d'une décharge, la structure de LixMn2O4 a été suivi par diffraction aux rayons X pour différentes compositions en lithium. L'évolution du paramètre $c$ de la maille hexagonale a été reportée ci-dessous.

\begin{minipage}{0.48\textwidth}
\includegraphics[width=\textwidth]{Ex12/Fig02}
\end{minipage}\hfill\begin{minipage}{0.48\textwidth}
\captionof{figure}{\'{E}volution du paramètre de maille en fonction de la st{\oe}chiométrie en lithium dans $\ce{Li_x Mn2O4}$.}
\end{minipage}

Conjointement, le titrage coulométrique a montré que pour une fraction molaire en dessous de $1$, \ce{Li_x Mn2O4} est constitué d'une seule phase cristalline.

\Q Qu'observe-t-on pour $x$ plus petit que $1$ ? Quelle hypothèse peut-on faire sur son origine ?
\solution{On observe que pour une st{\oe}chiométrie en lithium comprise entre 0.6 et 1.0, la structure cristallographique se dilate le long de l'axe $c$. Il est possible que cet allongement soit dûe à une déformation des octaèdres du réseau cristallin.}

Lors de l'acquisition de la force électromotrice pendant le titrage coulométrique, il a été remarqué que pour une st{\oe}chiométrie de 1, la f.e.m. s'écroule (il y a donc une perte de performance en puissance de la batterie).

\Q Comment évolue la fraction Mn(III) / Mn(IV) lors de la décharge ? Qu'en déduire ?
\solution{Lors de la décharge, on passe d'une structure contenant que des ions Mn(IV) à une structure contenant de manière st{\oe}chiométrique des ions Mn(III) et Mn(IV). \`{A} priori, l'augmentation du paramètre $c$ dans la structure cristallographique doit trouver son origine dans cette variation de charge.}

%\Q Donner la configuration électronique d'un ion \ce{Mn^{3+}}. Comment sont répartis les électrons dans les orbitales $3d$ ?


\Q En vous appuyant sur la théorie du champ cristallin, expliquer en quoi un octaèdre à la configuration électronique de \ce{Mn^{3+}} se distordra. Expliquer en quoi cette déformation garantit une stabilisation électronique de l'édifice moléculaire.
\solution{
Par exemple, si l'octaèdre \ce{MnO6} se déforme selon l'axe $z$, il s'en conduit une stabilisation des orbitales $d$ orientées selon cet axe : 
\begin{center}
\includegraphics[width=0.3\textwidth]{Ex12/Corr01}
\end{center}}

Cet effet, assez commun dans les métaux de transition est appelé effet Jahn-Teller. La distorsion de l'octaèdre se traduit macroscopiquement par une dilatation du matériau lors de la charge ou décharge.

\Q Expliquer pourquoi cette expansion est problématique dans le cas d'une batterie.
\solution{La dilatation pose problème puisqu'elle peut conduire à une surpression au sein des cellules de la batterie, et par conséquent, à sa dégradation.}


\Q Comment pouvons-nous remédier à ce problème ? 
\solution{On peut substituer l'élément manganèse par d'autres éléments Ni, Co dans lesquels l'effet Jahn-Teller n'est pas attendu.}


\paragraph*{Données}
\begin{itemize}
\item Numéro atomique : $Z ( \ce{Mn} ) = 25$
\item Degrés d'oxydation stable du manganèse : $+II$, $+III$, $+IV$, $+VI$ et $+VII$
\item Rayons ioniques : 
\begin{align*}
r( \ce{O^{2-}} ) &= 140 ~ \si{\pico\meter} ~ ; \\
r( \ce{Mn^{3+}}, \text{HS}) &= 65 ~ \si{\pico\meter} ~ ;  \\
r( \ce{Mn^{3+}}, \text{BS}) &= 58 ~ \si{\pico\meter} ~ ;  \\
r( \ce{Mn^{4+}} ) &= 53 ~ \si{\pico\meter}.
\end{align*}
\item \'{E}nergie de Stabilisation de Champ Cristallin de $Mn^{3+}$ pour une coordinence de [6] : 
\begin{align*}
\Delta_O = 2.60 ~ \si{\electronvolt}
\end{align*}
\end{itemize}
%r( \ce{Li^+} ) &= 76 ~ \si{\pico\meter} ; \\
%$r(O2-) = 140 pm ; r(Li+) = 76 pm ; r(Mn3+) = 67 pm$

\begin{comment}
\begin{center}
\begin{tabular}{ccc}
\toprule
& $d^3$ & $d^4$ \\
\midrule
octaédre $Oh$ HS & \multirow{2}{*}{$-6/5 \Delta_O$} & $-3/5 \Delta_O$ \\ 
octaédre $Oh$ BS & & $-8 /5 \Delta_O + P$ \\
tétraédre $Td$ & $-4/5 \Delta_T$ & $-2/5 \Delta_T$ \\
\bottomrule
\end{tabular}
\end{center}\end{comment}